\documentclass{report}
\usepackage{amsmath,amssymb}
\setlength{\parindent}{0mm}
\setlength{\parskip}{1em}
\begin{document}
\begin{center}
\rule{6in}{1pt} \
{\large Markus Huber \\
{\bf Resilience for Exascale Enabled Multigrid Methods}}

Insitute for Numerical Mathematics (M2) \\ Technische Universit�t M�nchen \\ Boltzmannstra�e 3 \\ 85748 Garching \\ Germany
\\
{\tt huber@ma.tum.de}\\
Bj�rn Gmeiner\\
Ulrich R�de\\
Barbara Wohlmuth\end{center}

With the increasing number of components and further miniaturization
the mean time between faults in supercomputers will decrease.
System level fault tolerance techniques
are expensive and cost energy, since they are often based on redundancy.
Also classical check-point-restart techniques reach their limits when the time
for storing the system state to backup memory becomes excessive.
Therefore, algorithm-based fault tolerance mechanisms can become an
attractive alternative.
This article investigates
the solution process for elliptic
partial differential equations that are discretized by
finite elements.
Faults that occur in the parallel geometric multigrid solver
are studied in various model scenarios.
In a standard domain partitioning approach,
the impact of a failure of a
core or a node will affect one or several subdomains.
Different strategies are developed to
compensate the effect of such a failure algorithmically.
The recovery is achieved by solving a
local subproblem with Dirichlet boundary conditions using
local multigrid cycling algorithms.
Additionally, we propose a superman strategy where extra
compute power is employed to minimize the time of the recovery process.


\end{document}
